\section{Functions available in models}

The functions documented in this section can be used in all mathematical expressions mentionned in Fig.~\ref{fig-syntax}.

The first two are general. The others are power system-oriented. Among them, some functions are restricted to some types of models.

\vspace*{2ex}

{\large \tt double precision function equal(var1,var2)}

\hspace*{10ex} {\tt double precision:: var1, var2}

{\it This function returns 1.d0 if $var1$ differs from $var2$ by less than $10^{-6}$, and 0.d0 otherwise.}

~ \hrulefill ~
%-----------------------------------------------------------------------

{\large \tt double precision function equalstr(var1,var2)}

\hspace*{10ex} {\tt character:: var1, var2}

{\it This function returns 1.d0 if the non-blank parts of str1 and str2 are the same,  and 0.d0  otherwise.}


~ \hrulefill ~
%-----------------------------------------------------------------------

{\large \tt double precision ppower([vx],[vy],[ix],[iy])}

{\it Returns the active power injected into network}

Can be used in models of type: {\tt inj}

Computation : 
\[ P = v_x i_x + v_y i_y \]

~ \hrulefill ~
%-----------------------------------------------------------------------
 
{\large \tt double precision qpower([vx],[vy],[ix],[iy])}

{\it Returns the reactive power injected into network}

Can be used in models of type: {\tt inj}

Computation : 
\[ Q = v_y i_x - v_x i_y \]
 
~ \hrulefill ~
%-----------------------------------------------------------------------

{\large \tt double precision function vrectif([if],[vin],\{kc\})}

\hspace*{10ex} {\tt double precision:: kc}

{\it This function is used to model the voltage drop in rectifiers. It gives the output voltage $vrectif$ as a function of the output current $if$ and the input voltage $vin$. It appears in some IEEE standard excitation models with $vrectif= f_{ex} \: E_{FD}$}

Can be used in models of type: {\tt exc}

Computation : 

\hspace*{20ex}\begin{minipage}[t]{15cm}
\begin{algorithmic} 
\STATE $in=kc \; if / max(vin, 10^{-3})$

\IF {$in \le 0$}
  \STATE $vrectif=vin$
\ELSIF {$in \le 0.433$}
  \STATE $vrectif=vin - 0.577 \; kc \;if$
\ELSIF {$in \le 0.75$}
  \STATE $vrectif= \sqrt{0.75 \; vin^2 - (kc \; if)^2}$
\ELSIF {$in \le 1.00$}
  \STATE $vrectif=1.732 (vin - kc \; if)$ 
\ELSE
  \STATE $vrectif=0$
\ENDIF
\end{algorithmic}
\end{minipage}

~ \hrulefill ~
%-----------------------------------------------------------------------

{\large \tt double precision function vinrectif([if],[vrectif],\{kc\})}

\hspace*{10ex} {\tt double precision:: kc}

{\it This function is the inverse of {\tt vrectif}. It is aimed at being called at initialization. For given values the field current $if$ and rectifier output voltage $vrectif$ it returns the value of the rectifier input voltage $vinrectif$. This determination is iterative because the portion of the nonlinear input-output characteristic on which the rectifier is operating is not known beforehand. The function does not work with a zero rectifier output $vrectif$ since, in this case, the input voltage is indeterminate.}

Can be used in models of type: {\tt exc}

Computation : 

\hspace*{20ex}\begin{minipage}[t]{15cm}
\begin{algorithmic}
\STATE $nbtries=1$
\STATE $vinest=vrectif+0.577 \; kc \;if$
\LOOP
\STATE $in=kc \; if / max(vinest,1d-03)$
\IF {$in \le 0$}
  \STATE $vinrectif=vrectif$
\ELSIF {$in \le 0.433$}
  \STATE $vinrectif=vrectif + 0.577 \; kc \;if$
\ELSIF {$in \le 0.75$}
  \STATE $vinrectif= \sqrt{vrectif^2 + (kc \; if)^2}/0.75$
\ELSE
  \STATE $vinrectif=(vrectif/1.732)+ kc \; if$
\ENDIF
\IF {$vinrectif= vinest$ or $nbtries>5$}
  \STATE exit loop
\ENDIF
\STATE $nbtries=nbtries+1$
\STATE $vinest=vinrectif$
\ENDLOOP
\end{algorithmic}
\end{minipage}

~ \hrulefill ~
%-----------------------------------------------------------------------

{\large \tt double precision function vcomp([v],[p],[q],\{Kv\},\{Rc\},\{Xc\})}

\hspace*{10ex} {\tt double precision:: Kv, Rc, Xc}

{\it This function returns the magnitude of a combination of the terminal voltage and the current of a synchronous machine. It appears in some IEEE standard excitation models}

Can be used in models of type: {\tt exc}

Computation :
\begin{eqnarray}
V_{comp} & = & | K_v \bar{V} + (R_c + j X_c) \bar{I} | = | K_v V + (R_c + j X_c) (i_P - j i_Q) | \nonumber \\
& = & | K_v V + (R_c + j X_c) (\frac{P}{V} - j \frac{Q}{V}) | \nonumber \\
& = & \frac{1}{V} \sqrt{(K_v V^2 + R_c P + X_c Q)^2 + (X_c P - R_c Q)^2} \nonumber
\end{eqnarray}

~ \hrulefill ~
%-----------------------------------------------------------------------

{\large \tt double precision function satur([ve],\{ve1\},\{se1\},\{ve2\},\{se2\})}

\hspace*{10ex} {\tt double precision:: ve1, se2, ve2, se2}

{\it This function returns the increment of field current needed to obtain a given output voltage, taking saturation into account. It appears in the model of some excitation systems.}

Can be used in models of type: {\tt exc}

Computation :
\[ satur= m \; ve^n \]
where:
\begin{eqnarray}
n & = & \frac{\log_{10}(se1/se2)}{\log_{10}(ve1/ve2)} \nonumber \\
m & = & \frac{se1}{ve1^n} \nonumber
\end{eqnarray}

{\bf Exception}. The function returns $satur=0$ if any of the following conditions holds true:
\[ ve \le 0 \mbox{~~~~or~~~~} ve1 \le 0 \mbox{~~~~or~~~~} ve2 \le 0 \mbox{~~~~or~~~~} ve1=ve2 \mbox{~~~~or~~~~} se1=0 \mbox{~~~~or~~~~} se2=0 \]

